\subsection{ソフトウェア保守の主要課題}
保守には、既存ソフトウェアを変更するという独自の難しさがある。保守担当者は、しばしば他人が開発した大規模で複雑なソフトウェアの中から原因を見つけ、変更し、予期しない影響を避ける必要がある。\par
\begin{itemize}
\item 技術的課題：限られた理解、テスト、影響分析、保守性。
\item 管理上の課題：組織目標との整合、人員配置、プロセス、供給者管理、組織構造。
\item 費用の課題：技術的負債と保守費用の見積り。
\item 測定の課題：どこに労力と費用が使われているかを見える化すること。
\end{itemize}
\subsubsection{技術的課題}
\paragraph{限られた理解}
限られた理解とは、保守担当者が、他者によって開発されたソフトウェアを最初は十分に理解していない状態である。変更すべき場所を見つけ、何を直せばよいかを理解するまでに大きな時間がかかる。\par
理解を難しくする要因には、領域知識の不足、読みにくい実装、文書不足、構成や提示の悪さがある。ソースコードのようなテキスト中心の表現では、変更履歴や設計意図を追うことが難しい。コード可視化、リバースエンジニアリング、明確な文書、コードブラウザが理解を助ける。\par
\paragraph{テスト}
保守では、変更要求や問題報告の処理中にテスト計画とテスト実行が行われる。大規模ソフトウェアで完全な再テストを繰り返す費用は高い。そのため、変更に応じて適切なテストを計画し、特に回帰テストを重視する必要がある。\par
回帰テストは、変更が意図しない副作用を起こしていないかを確認する選択的な再テストである。保守では複数の担当者が同時に異なる問題へ対応するため、テスト調整も難しくなる。さらに、重要機能を持つ本番システムを停止してテストすることが難しい場合もある。\par
\paragraph{影響分析}
\term{影響分析}{Impact Analysis}とは、提案された変更が既存ソフトウェアへ与える詳細な影響を評価する活動である。変更対象、影響を受ける構成要素、必要な資源、リスク、必要なテスト、性能・安全性・セキュリティへの影響を調べる。\par
影響分析では、変更要求や問題報告をまずソフトウェアの言葉へ翻訳する必要がある。たとえば「画面が遅い」という報告は、データベース照会、ネットワーク、キャッシュ、画面描画、外部サービスなど、複数の技術要素へ分解して考える必要がある。\par
\par\smallskip\noindent\refstepcounter{table}\label{tab:ch07-04}表 \thetable: 影響分析における影響分析で決めること・説明\par\smallskip
表 \thetable は、影響分析における影響分析で決めること・説明を比較・確認しやすい形で整理したものである。\par\smallskip
\begin{longtable}{p{0.30\linewidth}p{0.64\linewidth}}
\toprule
影響分析で決めること & 説明 \\
\midrule
\endhead
識別方法 & 変更要求や問題報告を一意に識別する方法を決める。 \\
分類と優先度 & 緊急度、重要度、保守分類、リスクで並べる。 \\
提出手順 & 利用者や運用者がどのように要求・報告を出すかを決める。 \\
追跡情報 & 利用者へ何を報告し、どの測定値で状況を示すかを決める。 \\
暫定回避策 & 本修正までの間、どのような回避策を提供するかを決める。 \\
フィードバック & 利用者へどの段階で何を返すかを決める。 \\
\bottomrule
\end{longtable}
\par\smallskip
\begin{center}
\centering
\IfFileExists{assets/generated_figures/ch07/fig-ch07-04.png}{\includegraphics[width=0.96\linewidth]{assets/generated_figures/ch07/fig-ch07-04.png}}{\fbox{\parbox[c][48mm][c]{0.9\linewidth}{\centering 画像生成待ち\\影響分析の連鎖図}}}
\par\smallskip
\refstepcounter{figure}\label{fig:ch07-04}図 \thefigure: 影響分析の連鎖図
\end{center}
図 \thefigure は、影響分析の連鎖図を視覚的に整理したものである。
\par\smallskip
\paragraph{保守性}
\term{保守性}{Maintainability}とは、ソフトウェア製品を修正できる能力である。修正には、欠陥の修正、改善、環境変化への適応、要求や機能仕様の変更が含まれる。保守性は重要な品質特性であり、開発中に仕様化、レビュー、管理されるべきである。\par
保守性が低いと、将来の保守担当者の負担が増え、技術的負債が蓄積する。緊急対応や短期的な機能追加で、十分なレビューや設計配慮を省くと、後からより多くの時間と労力が必要になる。\par
\begin{pointbox}{技術的負債を見る三つの観点}第一に、コード品質と業務上の重要度を分ける。すべての負債が同じ緊急度ではない。第二に、アーキテクチャが組織目標に合っているかを見る。第三に、チームやプロセス上の欠落が負債を増やしていないかを見る。\end{pointbox}
\subsubsection{管理上の課題}
\paragraph{組織目標との整合}
保守活動は、事業、顧客、利用者の優先順位と整合していなければならない。新規開発は、期限と予算を持つプロジェクトとして見えやすい。一方、保守は長期的に運用ソフトウェアの寿命を延ばし、価値を保つ活動であるため、経済効果が見えにくい。\par
その結果、リファクタリング、セキュリティ改善、性能改善よりも、目に見える新機能が優先されることがある。しかし、新機能だけを優先すると、コード劣化と技術的負債が増え、長期的な変更速度が落ちる。組織目標、保守性、工学標準の均衡が必要である。\par
\begin{itemize}
\item 運用中ソフトウェアの全ライフサイクル費用を理解する。
\item 新規開発と既存ソフトウェア改善の費用対効果を比較する。
\item 同じチームが開発と保守に責任を持つことで、技能と責任の分断を減らす。
\item 初期開発から保守性要求に注意を向ける。
\end{itemize}
\paragraph{人員配置}
保守作業は、しばしば新規開発より魅力が低いと見られる。しかし、保守はライフサイクル活動の大きな部分を占め、保守担当者の貢献は重要である。組織は、保守担当者の専門性を評価し、学習機会とキャリア上の意味を与える必要がある。\par
人員配置では、既存コードを読み解く能力、運用への理解、テスト設計、障害対応、コミュニケーション能力が必要である。保守担当者が孤立すると、知識が属人化し、対応品質が下がる。\par
\paragraph{プロセス}
ソフトウェアライフサイクルプロセスは、ソフトウェアと関連成果物を開発・保守するための活動、方法、実践、変換の集合である。保守は開発と多くを共有する。たとえば、構成管理、変更管理、レビュー、テスト、文書更新はどちらにも必要である。\par
一方で、保守には開発にない活動もある。既存プログラム理解、運用からの問題報告処理、ヘルプデスク連携、緊急修正、リリース後の小さな変更群の管理がその例である。\par
\paragraph{供給者管理}
供給者管理とは、保守を外部供給者に委託する場合に、供給者とその成果を適切に管理することである。外部委託、オフショア保守、複数供給者によるサービス提供は、現代の保守では一般的な選択肢である。\par
\par\smallskip\noindent\refstepcounter{table}\label{tab:ch07-05}表 \thetable: 供給者管理における契約形態・説明\par\smallskip
表 \thetable は、供給者管理における契約形態・説明を比較・確認しやすい形で整理したものである。\par\smallskip
\begin{longtable}{p{0.30\linewidth}p{0.64\linewidth}}
\toprule
契約形態 & 説明 \\
\midrule
\endhead
単一供給者・提携 & 一つの供給者がすべての保守サービスを提供する。または外部の統合者が複数供給者との関係を管理する。 \\
複数供給者 & 複数の独立した供給者が製品やサービスを提供し、それらを単一のソフトウェア支援サービスとして統合する。XaaS の普及により増えている。 \\
\bottomrule
\end{longtable}
外部委託では、保守範囲、サービス水準合意、契約上の責任、連絡体制、ヘルプデスク、プロセス自動化を明確にする必要がある。特に中核業務を支えるソフトウェアでは、組織が制御を失わないように注意する。\par
\paragraph{保守の組織面}
保守組織の設計では、誰が保守に責任を持つかを決める。アジャイル型では、開発者が保守も担い、開発者兼保守担当者として働くことが多い。一方で、開発チームとは別の保守チームを置く組織もある。\par
\par\smallskip\noindent\refstepcounter{table}\label{tab:ch07-06}表 \thetable: 保守の組織面における方式・利点・注意点\par\smallskip
表 \thetable は、保守の組織面における方式・利点・注意点を比較・確認しやすい形で整理したものである。\par\smallskip
\begin{longtable}{p{0.22\linewidth}p{0.40\linewidth}p{0.30\linewidth}}
\toprule
方式 & 利点 & 注意点 \\
\midrule
\endhead
同一チーム方式 & 開発知識が残り、要求変化へ素早く対応しやすい。 & 障害対応が新規開発を中断しやすい。文書化が軽視されることがある。 \\
分離チーム方式 & 運用と保守に集中でき、長期支援を組織化しやすい。 & 引き継ぎ摩擦、知識喪失、保守担当者の動機低下が起こり得る。 \\
\bottomrule
\end{longtable}
\begin{pointbox}{判断基準}どの組織形態が正しいかは一律には決まらない。重要なのは、経験ある人やチームへ保守責任を委ね、保守作業と変更履歴に関する品質の高い文書を残すことである。\end{pointbox}
\subsubsection{保守費用}
保守費用を理解するには、保守分類ごとに費用傾向を見る必要がある。是正、予防、適応、追加、完全化のどこに努力が使われているかを示せば、顧客や組織は保守の価値と課題を理解しやすくなる。\par
\paragraph{技術的負債の費用見積り}
技術的負債は、コードを必要以上に保守しにくくし、変更費用を増やす。技術的負債の費用を見積もるには、規模、複雑さ、設計原則違反、コーディング慣行違反、コードのにおい、文書不足、テスト不足などを見る。\par
保守性を改善するには、読みやすさを確保し、構造化されたコードを保ち、複雑さを減らし、正確なコメントを付け、言語上の弱点やコンパイラ依存構造を避け、エラー追跡を容易にし、設計とコードの追跡可能性を保ち、コードレビューを行う必要がある。\par
\par\smallskip
\begin{center}
\centering
\IfFileExists{assets/generated_figures/ch07/fig-ch07-05.png}{\includegraphics[width=0.96\linewidth]{assets/generated_figures/ch07/fig-ch07-05.png}}{\fbox{\parbox[c][48mm][c]{0.9\linewidth}{\centering 画像生成待ち\\技術的負債と保守費用の関係図}}}
\par\smallskip
\refstepcounter{figure}\label{fig:ch07-05}図 \thefigure: 技術的負債と保守費用の関係図
\end{center}
図 \thefigure は、技術的負債と保守費用の関係図を視覚的に整理したものである。
\par\smallskip
\paragraph{保守費用見積り}
保守費用の見積りは、ソフトウェア計画の早い段階で作成すべきである。見積りは保守活動の範囲に依存する。保守が進むにつれて、過去の測定データを使って見積りを更新する。\par
\par\smallskip\noindent\refstepcounter{table}\label{tab:ch07-07}表 \thetable: 保守費用見積りにおける費用要素・説明\par\smallskip
表 \thetable は、保守費用見積りにおける費用要素・説明を比較・確認しやすい形で整理したものである。\par\smallskip
\begin{longtable}{p{0.30\linewidth}p{0.64\linewidth}}
\toprule
費用要素 & 説明 \\
\midrule
\endhead
利用者拠点への移動 & 現地支援や調査が必要な場合の費用である。 \\
訓練 & 保守担当者と利用者への教育費用である。 \\
開発・テスト環境 & ソフトウェア工学環境、テスト環境、ツールの維持費である。 \\
人件費 & 給与、福利厚生、要員確保にかかる費用である。 \\
消耗品・設備 & 保守作業に必要な追加資源である。 \\
ライセンス保守 & 商用ソフトウェアや基盤ライセンスの維持費である。 \\
製品変更と管理 & リリース管理、プログラム管理、変更実施にかかる費用である。 \\
\bottomrule
\end{longtable}
個別の変更要求や問題報告についても、影響分析の中で費用見積りが必要である。見積りには、実施工数、必要資源、実施期間、リスクを含める。パラメトリックモデル、類似システム比較、経験データなどの方法を使う。\par
\subsubsection{ソフトウェア保守の測定}
測定可能な保守対象には、保守プロセス、資源、製品がある。測定値には、規模、複雑さ、品質、理解しやすさ、保守性、労力が含まれる。特に、保守分類ごとに使われた労力を測ると、費用構造が見える。\par
\par\smallskip\noindent\refstepcounter{table}\label{tab:ch07-08}表 \thetable: ソフトウェア保守の測定における測定対象・説明\par\smallskip
表 \thetable は、ソフトウェア保守の測定における測定対象・説明を比較・確認しやすい形で整理したものである。\par\smallskip
\begin{longtable}{p{0.30\linewidth}p{0.64\linewidth}}
\toprule
測定対象 & 説明 \\
\midrule
\endhead
保守性 & モジュール性、再利用性、解析容易性、修正容易性、テスト容易性、支援容易性などで測る。 \\
信頼性 & 成熟性、可用性、耐故障性、回復性などで測る。 \\
規模 & 機能規模、コード行数、構成要素数などで測る。 \\
要求数 & 期間ごとの変更要求数、問題報告数、分類別件数を見る。 \\
労力 & 一件あたりの保守工数、分類別工数、リリースあたり工数を見る。 \\
技術特性 & プラットフォーム、言語、フレームワーク、ハードウェアなどを記録する。 \\
\bottomrule
\end{longtable}
\begin{pointbox}{測定の目的}測定は管理のためだけではない。保守費用がどこに使われているかを説明し、改善対象を見つけ、組織内でアプリケーションごとの保守状況を比較するために使う。\end{pointbox}
